% !TEX TS-program = pdflatex
% !TeX program = pdflatex
% !TEX encoding = UTF-8
% !TEX spellcheck = fr

\documentclass[xcolor=table]{beamer}

\title[Thèse : Vers une amélioration des RAT] %
{Vers une amélioration \\des résumés automatiques de textes} 

\institute{ %
%	\begin{flushleft}
		École nationale Supérieure d'Informatique (ESI, ex. INI), Alger, Algérie 
		
		École doctorale STIC (2010/2011)
		
		Spécialité : Informatique répartie et mobile (IRM)
		
		Laboratoire LCSI - Groupe D3
		
		Année d'inscription en doctorat : (2013/2014)	
%	\end{flushleft}
}

\author[ \textbf{\footnotesize\insertframenumber/\inserttotalframenumber} \hspace*{1.5cm} ESI - ARIES (2020)] %
{\Large\bfseries ARIES Abdelkrime}

\titlegraphic{%
	\includegraphics[height=1cm]{img/esi-logo.png} 
	\hfill
	28 Octobre 2020
	\hfill
	\includegraphics[height=1cm]{img/K-LCSI.png}
}%\hspace*{4.75cm}~

\date{\footnotesize
Directeur de thèse : \textbf{Pr. ZEGOUR Djamal Eddine} \\
Co-directeur de thèse : \textbf{Pr. HIDOUCI Walid Khaled}
} %\today

\usetheme{Warsaw} % Antibes Boadilla Warsaw

\beamertemplatenavigationsymbolsempty



%\usepackage{fullpage}
%\usepackage[left=2.8cm,right=2.2cm,top=2 cm,bottom=2 cm]{geometry}
\setbeamersize{text margin left=10pt,text margin right=10pt}
\usepackage{amsmath,amssymb} 
\usepackage[T1]{fontenc}
\usepackage{textcomp}
\usepackage[utf8]{inputenc}
\usepackage[french, onelanguage]{algorithm2e}
\usepackage[french]{babel}
\usepackage{arabtex}
\usepackage{txfonts}
\usepackage[]{graphicx}
\usepackage{multirow}
\usepackage{hyperref}
\usepackage{colortbl}
\usepackage{listingsutf8}
\usepackage{wrapfig}
\usepackage{multicol}
\usepackage[export]{adjustbox} %for images in table, also for frame

\hypersetup{
	colorlinks,
	urlcolor = blue
}

%\renewcommand{\baselinestretch}{1.5}

\def\supit#1{\raisebox{0.8ex}{\small\it #1}\hspace{0.05em}}

\AtBeginSection{%
	\begin{frame}
		\sectionpage
	\end{frame}
}

\newcommand{\rottext}[2]{%
	\rotatebox{90}{%
	\begin{minipage}{#1}%
		\raggedleft#2%
	\end{minipage}%
	}%
}

\usepackage{longtable}
\usepackage{tabu}


%\setbeamertemplate{headline}{}

\definecolor{lightblue}{HTML}{D0D2FF}
\definecolor{lightyellow}{HTML}{FFFFAA}
\definecolor{darkblue}{HTML}{0000BB}
\definecolor{olivegreen}{HTML}{006600}
\definecolor{violet}{HTML}{6600CC}

\newcommand{\keyword}[1]{\textcolor{red}{\bfseries\itshape #1}}
\newcommand{\expword}[1]{\textcolor{olivegreen}{#1}}
\newcommand{\optword}[1]{\textcolor{violet}{\bfseries #1}}

\makeatletter
\newcommand\mysphere{%
	\parbox[t]{10pt}{\raisebox{0.2pt}{\beamer@usesphere{item projected}{bigsphere}}}}
\makeatother

%\let\oldtabular\tabular
%\let\endoldtabular\endtabular
%\renewenvironment{tabular}{\rowcolors{2}{white}{lightblue}\oldtabular\rowcolor{blue}}{\endoldtabular}


\NoAutoSpacing %french autospacing after ":"

\def\graphpath{}

\newcommand{\changegraphpath}[1]{\def\graphpath{#1}}


\newcommand{\vgraphpage}[2][.84\textheight]{%
%	\begin{center}%
		\includegraphics[height=#1]{\graphpath #2}%
%	\end{center}%
}

\newcommand{\hgraphpage}[2][\textwidth]{%
%	\begin{center}%
		\includegraphics[width=#1]{\graphpath #2}%
%	\end{center}%
}

\newcommand{\graphpage}[2][]{%
	\includegraphics[#1]{\graphpath #2}%
}

\bibliographystyle{apalike}

\newcommand{\insertbibliography}[2]{
	\appendix
	\section*{Bibliographie}
	\nocite{#2}
%	\makeatletter % to change template
%	\setbeamertemplate{headline}[default] % not mandatory, but I though it was better to set it blank
%	\def\beamer@entrycode{\vspace*{-\headheight}} % here is the part we are interested in :)
%	\makeatother
	\begin{multicols*}{2}[\frametitle{\insertsection} \usebeamertemplate*{frametitle}]%\usebeamertemplate*{frametitle}\frametitle{Références}
		\tiny
		\bibliography{#1}
	\end{multicols*}
}

\newcommand{\insertlicence}{
	\begin{frame}[plain]
	\frametitle{Licence : CC-BY-NC 4.0}
%	\framesubtitle{Licence: CC-BY-NC 4.0}
	\begin{center}
		\bfseries Attribution - Pas d'Utilisation Commerciale 4.0 International 
		\url{https://creativecommons.org/licenses/by-nc/4.0/deed.fr}
	\end{center}

	\tiny

	Vous êtes autorisé à : 
	\begin{itemize}
		\item \textbf{Partager} — copier, distribuer et communiquer le matériel par tous moyens et sous tous formats
		\item \textbf{Adapter} — remixer, transformer et créer à partir du matériel
	\end{itemize}
	
	Selon les conditions suivantes : 
	\begin{itemize}
		\item \textbf{Attribution} — Vous devez créditer l'Œuvre, intégrer un lien vers la licence et indiquer si des modifications ont été effectuées à l'Oeuvre. Vous devez indiquer ces informations par tous les moyens raisonnables, sans toutefois suggérer que l'Offrant vous soutient ou soutient la façon dont vous avez utilisé son Oeuvre.
		\item \textbf{Pas d'Utilisation Commerciale} — Vous n'êtes pas autorisé à faire un usage commercial de cette Oeuvre, tout ou partie du matériel la composant. 
		\item \textbf{Pas de restrictions complémentaires} — Vous n'êtes pas autorisé à appliquer des conditions légales ou des mesures techniques qui restreindraient légalement autrui à utiliser l'Oeuvre dans les conditions décrites par la licence.
	\end{itemize}

	\end{frame}
}

\settowidth{\leftmargini}{\usebeamertemplate{itemize item}}
\addtolength{\leftmargini}{\labelsep}

\AtBeginDocument{
	\newcolumntype{L}[2]{>{\vbox to #2\bgroup\vfill\flushleft}p{#1}<{\egroup}} 
	
	\begin{frame}[plain]
		\maketitle
	\end{frame}

%	\insertlicence
}


% needs etoolbox; to break links after -
\appto\UrlBreaks{\do\-}

\definecolor{darkgreen}{RGB}{18,118,34}
\definecolor{darkred}{RGB}{201,33,30}
\definecolor{darkblue}{RGB}{0,0,139}
\definecolor{darkyellow}{RGB}{150,150,0} 

\begin{document}

\begin{frame}
\frametitle{Résumé automatique de textes}
\framesubtitle{Motivation}

\vspace{-12pt}
\hgraphpage{img/motivation.pdf}

%\begin{center}
%	\hgraphpage[0.5\textwidth]{img/motivation.pdf}
%\end{center}
%
%\begin{itemize}
%	\item En plus, les employés peuvent lire un résumé des messages du directeur de l'entreprise, les étudiants peuvent lire l'essentiel de ce que les enseignants envoient, etc.
%\end{itemize}

\end{frame}

\begin{frame}
	\frametitle{Résumé automatique de textes : Introduction}
	\framesubtitle{Classification du résumé automatique}
%	\vspace{-12pt}
	\hgraphpage{img/sum-classif.pdf}
	
\end{frame}

\begin{frame}
	\frametitle{Résumé automatique de textes : Introduction}
	\framesubtitle{Problématique}
	
	\vspace{-6pt}
	\begin{itemize}
		\item \optword{Définition} : idées plus importantes, texte plus concise. 
		\begin{itemize}
			\item Comment savoir ce qu'il est important ?
			\item Par rapport à qui/quoi ?
		\end{itemize}
		\item \optword{Informativité} : un résumé représentatif  
		\begin{itemize}
			\item Couverture : le résumé couvre-t-il tout le contenu original ?
			\item Concision : y a t-il de redondance dans le résumé ?
		\end{itemize}
		\item \optword{Lisibilité} : un résumé proche de celui des humains
		\begin{itemize}
			\item Clarté des références : les références sont-elles claires ?
			\item Cohérence : les phrases ont-elles un ordre logique ?
			\item Grammaticalité : les phrases sont-elles grammaticalement correctes ?
		\end{itemize}
		\item \optword{Ressources} : implémentation facile de la méthode
		\begin{itemize}
			\item Outils : existent-ils des outils pour la langue destinataire ?
			\item Corpus : existent-ils des corpus pour la langue destinataire ?
		\end{itemize}
		\item \optword{Évaluation} : juger les aspects précédents dans un résumé 
		\begin{itemize}
			\item Intrinsèque vs. Extrinsèque
			\item Manuelle vs. Automatique
		\end{itemize}
	\end{itemize}
	
\end{frame}

\begin{frame}
\frametitle{Résumé automatique de textes : Introduction}
\framesubtitle{Objectifs}
	
	\begin{itemize}
		\item \optword{Adaptation facile} : utilisation pour différents domaines et langues 
		\begin{itemize}
			\item Méthode : méthode indépendante de langue et du domaine
			\item Implémentation : système facile à mettre à jour
		\end{itemize}
		\item \optword{Résumés informatifs} : décrivent le texte original
		\begin{itemize}
			\item Couverture : plus d'information utile
			\item Concision : moins de redondance
		\end{itemize}
		\item \optword{Cohérence} : Ordre des phrases
		\item \optword{Rapidité} : génération rapide des résumés
%		\begin{itemize}
%			\item Couverture : plus d'information utile
%			\item Concision : moins de redondance
%		\end{itemize}
	\end{itemize}
	
\end{frame}

\begin{frame}
\frametitle{Résumé automatique de textes}
\framesubtitle{Plan}

\begin{multicols}{2}
	%	\small
	\tableofcontents
\end{multicols}
\end{frame}

% ================================================
%
% ================================================
\section{Approches de résumé automatique} 

\begin{frame}
\frametitle{Résumé automatique de textes}
\framesubtitle{Approches de résumé automatique}

\begin{columns}
\begin{column}{0.32\textwidth}
	\begin{block}{\scriptsize\bfseries\cite{12-nenkova-mckeown}}
		\begin{itemize}
			\item représentation du sujet 
			:
			mots du sujet,
			fréquences, 
			analyse sémantique latente, 
			modèles de sujets bayésiens,
			clustering
%			\begin{itemize}
%				\item mots du sujet 
%				\item fréquences
%				\item analyse sémantique latente
%				\item modèles de sujets bayésiens
%			\end{itemize}
			\item représentation des indicateurs : 
			par graphes, 
			apprentissage automatique
		\end{itemize}
	\end{block}
\end{column}
\begin{column}{0.3\textwidth}
	\begin{block}{\scriptsize\bfseries\cite{12-lloret-palomar}}
		\begin{itemize}
			\item statistique 
			\item par graphes
			\item basée discours
			\item par apprentissage automatique
		\end{itemize}
	\end{block}
\end{column}
\begin{column}{0.28\textwidth}
	\begin{block}{\scriptsize\bfseries\cite{19-aries-al}}
		\begin{itemize}
			\item statistique 
			\item par graphes
			\item linguistique 
			\item par apprentissage automatique
		\end{itemize}
	\end{block}
\end{column}
\end{columns}
	
\end{frame}

\subsection{Approche statistique} 

\begin{frame}
\frametitle{Résumé automatique de textes : Approches}
\framesubtitle{Approche statistique}

\begin{itemize}
	\item \optword{Méthodes}
	\vspace{-.5cm}
	\begin{multicols}{3}
		\begin{itemize}
			\item Fréquence
			\item Position
			\item Mots de titre
			\item Longueur de la phrase
			\item Centroid
			\item Frequent itemsets
			\item Analyse sémantique latente
		\end{itemize}
	\end{multicols}

	\item \optword{Avantages}
	\vspace{-.5cm}
	\begin{multicols}{2}
		\begin{itemize}
			\item Moins de ressources
			\item Simple et rapide
		\end{itemize}
	\end{multicols}

	\item \optword{Inconvénients}
	\vspace{-.5cm}
	\begin{multicols}{2}
		\begin{itemize}
			\item Lisibilité
			\item Seulement extractive
		\end{itemize}
	\end{multicols}
	
	\item \optword{Problèmes et solutions probables}
	\begin{itemize}
		\item combinaison des caractéristiques : Apprentissage automatique
		\item pertinence : caractéristiques linguistiques
		\item redondance dans les phrases : compression
	\end{itemize}

\end{itemize}
	
\end{frame}

\subsection{Approche par graphes} 

\begin{frame}
\frametitle{Résumé automatique de textes : Approches}
\framesubtitle{Approche par graphes}

\begin{itemize}
	\item \optword{Méthodes}
	\vspace{-.5cm}
	\begin{multicols}{2}
		\begin{itemize}
			\item Propriétés de graphe
			\item Itérative
		\end{itemize}
	\end{multicols}
	
	\item \optword{Avantages}
	\vspace{-.5cm}
	\begin{multicols}{3}
		\begin{itemize}
			\item Moins de ressources
			\item Simple
			\item Cohérence
		\end{itemize}
	\end{multicols}
	
	\item \optword{Inconvénients}
	\vspace{-.5cm}
	\begin{multicols}{2}
		\begin{itemize}
			\item Puissance de calcul
		\end{itemize}
	\end{multicols}
	
	\item \optword{Problèmes et solutions probables}
	\begin{itemize}
		\item similarité entre phrases : TF-IDF, sémantique
		\item propriétés de la phrase : caractéristiques statistiques et linguistiques, similarité phrase/document, propriété temporelle
	\end{itemize}
	
\end{itemize}
	
\end{frame}

\subsection{Approche linguistique} 

\begin{frame}
\frametitle{Résumé automatique de textes : Approches}
\framesubtitle{Approche linguistique}
	
\begin{itemize}
	\item \optword{Méthodes}
	\vspace{-.5cm}
	\begin{multicols}{3}
		\begin{itemize}
			\item mots de sujet
			\item indicateurs
			\item co-occurrence
			\item structure rhétorique
		\end{itemize}
	\end{multicols}
	
	\item \optword{Avantages}
	\vspace{-.5cm}
	\begin{multicols}{3}
		\begin{itemize}
			\item plus de précision
			\item abstractif
		\end{itemize}
	\end{multicols}
	
	\item \optword{Inconvénients}
	\vspace{-.5cm}
	\begin{multicols}{3}
		\begin{itemize}
			\item ressources (outils)
			\item complexe
			\item puissance de calcul
		\end{itemize}
	\end{multicols}
	
	\item \optword{Problèmes et solutions probables}
	\begin{itemize}
		\item règles de génération : Apprentissage automatique
	\end{itemize}
	
\end{itemize}

\end{frame}

\subsection{Approche par apprentissage automatique} 

\begin{frame}
\frametitle{Résumé automatique de textes : Approches}
\framesubtitle{Approche par apprentissage automatique}
	
\begin{itemize}
	\item \optword{Méthodes}
	\vspace{-.5cm}
	\begin{multicols}{3}
		\begin{itemize}
			\item fonction de réglage
			\item fonction de décision
			\item modèles de sujets bayésiens
			\item apprentissage par renforcement
			\item apprentissage profond
		\end{itemize}
	\end{multicols}
	
	\item \optword{Avantages}
	\vspace{-.5cm}
	\begin{multicols}{1}
		\begin{itemize}
			\item estimation automatique de règles
		\end{itemize}
	\end{multicols}
	
	\item \optword{Inconvénients}
	\vspace{-.5cm}
	\begin{multicols}{1}
		\begin{itemize}
			\item ressources (corpus)
		\end{itemize}
	\end{multicols}
	
	\item \optword{Problèmes et solutions probables}
	\begin{itemize}
		\item données étiquetées : Apprentissage par renforcement
		\item insuffisance de données : création de corpus
	\end{itemize}
	
\end{itemize}

\end{frame}


% ================================================
%
% ================================================
\section{Baseline (AllSummarizer-TCC)} 

\begin{frame}
\frametitle{Résumé automatique de textes}
\framesubtitle{Baseline (AllSummarizer-TCC)}
	
	\hgraphpage{img/tcc-arch.pdf}
	
\end{frame}

\subsection{Prétraitement}

\begin{frame}
\frametitle{Résumé automatique de textes : AllSummarizer-TCC}
\framesubtitle{Prétraitement}
	
\begin{block}{Outils utilisés : tâche/langues}
	\scriptsize\bfseries
	\begin{tabular}{p{.2\textwidth}p{.2\textwidth}p{.5\textwidth}} 
		\hline \hline
		Tâche & Outils & Langues \\
		\hline
		\multirow{3}{2cm}{Segmentation de phrases} & openNLP & Nl, En, De, It, Pt, Th \\
		\cline{2-3}
		& JHazm & Fa \\
		\cline{2-3}
		& Regex & Le reste \\
		\hline
		\multirow{3}{2cm}{Séparation de mots} & openNlp & Nl, En, De, It, Pt, Th \\
		\cline{2-3}
		& Lucene & Zh, Ja \\
		\cline{2-3}
		& Regex & Le reste \\
		\hline
		\multirow{3}{2cm}{Racinisation} & Shereen Khoja & Ar \\
		\cline{2-3}
		& JHazm & Fa \\
		\cline{2-3}
		& HebMorph & He \\
		\cline{2-3}
		& Lucene & Bg, Cs, El, Hi, Id, Ja, No \\
		\cline{2-3}
		& Snowball & Eu, Ca, Nl, En (Porter), Fi, Fr, De, Hu, It, Pt, Ro, Ru, Es, Sv, Tr \\
		\cline{2-3}
		& / & Le reste \\
		\hline \hline
	\end{tabular}
\end{block}

\end{frame}

\subsection{Traitement}

\begin{frame}
\frametitle{Résumé automatique de textes : AllSummarizer-TCC}
\framesubtitle{Traitement}
	
	\begin{itemize}
		\item Un texte peut contenir plusieurs sujets, et une phrase peut discuter plusieurs sujets
		\begin{itemize}
			\item Regroupement de phrases avec similarité et seuil de regroupement ($Th$)
		\end{itemize}
		\item Une phrase est importante si elle peut représenter le maximum des sujets(cluters)
		\[ Score(s_i , c_j , f_k ) = 1 + \sum_{\phi \in s_i} {P(f_k=\phi | s_i \in c_j)} \]
		\[ Score(s_i , \bigcap_{j} c_j , F) =  %\propto 
		\prod_{j} \prod_{k} Score(s_i , c_j , f_k ) \]
		$ s $ : phrase, $ c $ : cluster, $ f $ : caractéristique, $ F $ : ensemble de caractéristiques, $ \phi $: observation de $ f $.
		\item $f$ : TF (Uni, Bi); Pos (intervalle de 10); Len (réel, pré-traité)
	\end{itemize}
	
\end{frame}

\subsection{Extraction}

\begin{frame}
\frametitle{Résumé automatique de textes : AllSummarizer-TCC}
\framesubtitle{Extraction}
	
\hgraphpage{img/tcc-extract.pdf}
	
\end{frame}

\subsection{Expérimentation} 

\begin{frame}
\frametitle{Résumé automatique de textes : Baseline}
\framesubtitle{Expérimentation}
	
	\begin{itemize}
		\item Participation dans le workshop MultiLing'15 \cite{15-aries-al}
		\item Procédure :
		\begin{itemize}
			\item Un dataset d'entraînement a été envoyé aux participants (avec les résumés de référence)
			\item Les systèmes peuvent être entraînés/réglés en utilisant ce dataset
			\item Un dataset de test a été envoyé aux participants (sans résumés de référence)
			\item Les participants doivent générer des résumés et les envoyer
			\item La comité d'organisation s'occupe de l'évaluation des résultats et la diffusion des résultats
		\end{itemize}
		\item Tâches (dans lesquelles nous avons participé) :
		\begin{itemize}
			\item MSS : résumé multilingues mono-document
			\item MMS : résumé multilingues multi-documents
		\end{itemize}
	\end{itemize}
	
\end{frame}

\begin{frame}
\frametitle{Résumé automatique de textes : Baseline}
\framesubtitle{Expérimentation : Réglage de paramètres}
	
\begin{itemize}
	
	\item Estimation de seuil de clustering 
	\vspace{-.5cm}
	\begin{multicols}{3}
	\begin{itemize}
		\item Médiane
		\item Moyenne
		\item Mode 
		\item Variance
		\item $ sDn = \frac{\sum |s|}{|D| * n}$ 
		\item $ Dsn = \frac{|D|}{n * \sum |s|}$ 
		\item $ Ds = \frac{|D|}{\sum |s|}$ 
	\end{itemize}
	\end{multicols}
	\vspace{-.4cm}
	$|s|$: number of different terms in a sentence $s$. 
	$|D|$: number of different terms in the document $D$.
	$n$: number of sentences in this document. 
	
	\item Sélection des caractéristiques 
	\begin{itemize}
		\item Pour chaque langue, trouver les combinaisons de caractéristiques qui maximisent ROUGE-2 dans chaque tâche (MSS et MMS)
	\end{itemize}
\end{itemize}
	
\end{frame}

\begin{frame}
\frametitle{Résumé automatique de textes : Baseline}
\framesubtitle{Expérimentation : Évaluation}
	
TCC = Notre système (AllSummarizer-TCC)\\
S = un système participant avec n langues\\

\[ 
AVG_{S} = \frac{\sum\limits_{i=1}^{n} Score_S(L_i)}{n}
\hskip0.2\textwidth 
RI = \frac{AVG_{TCC} - AVG_{S}}{AVG_{S}}
\]

Métriques : 
\begin{itemize}
	\item ROUGE \cite{04-lin}
	\item AutoSummENG et MeMoG \cite{11-giannakopoulos-al}
	\item NPowER \cite{13-giannakopoulos-karkaletsis}
\end{itemize}
	
\end{frame}

\begin{frame}
\frametitle{Résumé automatique de textes : Baseline}
\framesubtitle{Expérimentation : Résultats MSS}
	
\begin{block}{Amélioration relative de TCC dans MSS}
	\scriptsize\bfseries
	\begin{tabular}{p{.28\textwidth}p{.1\textwidth}p{.1\textwidth}p{.1\textwidth}p{.1\textwidth}p{.1\textwidth}} 
		\hline \hline
		\multirow{2}{*}{Methods} & \multicolumn{5}{c}{Amélioration relative \%}\\
		\cline{2-6}
		& R-1	& R-2	& R-3	& R-4	& R-SU4\\
		\hline
		BGU-SCE-M (ar, en, he)		& -09.19	& -14.02	& -19.39	& -25.12	& -11.07\\
		EXB (all 38)				& -07.64	& -10.55	& -09.86	& -07.92	& -10.63\\
		CCS (all 38)				& -07.33	& -13.24	& -10.95	& -03.04	& -07.40\\
		BGU-SCE-P (ar, en, he)		& -04.33	& -01.63	& -02.69	& -06.16	& -01.89\\
		UA-DLSI (en, de, es)		& +02.12	& +06.25	& +13.86	& +17.15	& +05.62\\
		NTNU (en, zh)				& +06.44	& +07.06	& +11.50	& +21.81	& +05.74\\
		\hline
		Oracle (all 38) [TopLine]	& -31.64	& -49.00	& -63.80	& -72.91	& -36.77\\
		Lead (all 38) [BaseLine]	& +02.39	& +08.67	& +08.20	& +04.02	& +05.82\\
		\hline \hline
	\end{tabular}
\end{block}
	
\end{frame}

\begin{frame}
\frametitle{Résumé automatique de textes : Baseline}
\framesubtitle{Expérimentation : Résultats MSS}
	
\begin{block}{Amélioration relative de TCC dans MMS}
	\scriptsize\bfseries
	\begin{tabular}{p{.4\textwidth}p{.15\textwidth}p{.15\textwidth}p{.15\textwidth}} 
		\hline \hline
		\multirow{2}{*}{SysID} & \multicolumn{3}{c}{Our method improvement \%}\\
		\cline{2-4}
		& AutoSummENG 	& MeMoG 	& NPowER \\
		\hline
		UJF-Grenoble (fr, en, el) 	& -08.87 		& -14.55	& -03.62 \\
		UWB (all 10) 		& -22.56 		& -22.66 	& -07.54 \\
		ExB (all 10) 		& -09.44 		& -09.16 	& -02.80 \\
		IDA-OCCAMS (all 10) 		& -17.11 		& -17.68 	& -05.53 \\
		GiauUngVan (- zh, ro, es)	& -16.43		& -19.40	& -05.68 \\
		SCE-Poly (ar, en, he)	& -05.72		& -03.35	& -01.46 \\
		BUPT-CIST (all 10)		& +10.67		& +11.53	& +02.85 \\
		BGU-MUSE (ar, en ,he)	& +05.67		& +06.92	& +01.74 \\
		NCSR/SCIFY-NewSumRerank (- zh)		& +01.53		& -01.25	& +00.13 \\
		\hline
		AllSummazer (MSS param) (all 10)		& +01.98		& +02.35	& +00.58 \\
		\hline \hline
	\end{tabular}
\end{block}
	
\end{frame}

% ================================================
%
% ================================================
\section{Scores cumulatifs à base de graphes} 

\begin{frame}
\frametitle{Résumé automatique de textes}
\framesubtitle{Scores cumulatifs à base de graphes}
	
	\begin{center}
		\vgraphpage{img/gc-archi.pdf}
	\end{center}
	
\end{frame}

\subsection{Génération des phrases candidates}

\begin{frame}
\frametitle{Résumé automatique de textes : Scores cumulatifs}
\framesubtitle{Génération des phrases candidates : Génération du graphe}
	
\begin{columns}
	\begin{column}{0.45\textwidth}
		\begin{itemize}
			\item Les phrases sont pré-traitées
			\item Un graphe $ G(V, E) $ est construit 
			\begin{itemize}
				\item $V$ : les phrases comme nœuds
				\item $E$ : les similarités entre phrases comme arcs
			\end{itemize}
			\item Similarité : cosinus 
			\[cos(X,Y) = \frac {\sum_i {x_i.y_i} }
			{\sqrt{\sum_i(x_i)^2} . \sqrt{\sum_i(y_i)^2}}\]
		\end{itemize}
	\end{column}
	\begin{column}{0.45\textwidth}
		\hgraphpage{img/sim-graphe.pdf}
	\end{column}
\end{columns}

\end{frame}

\begin{frame}
\frametitle{Résumé automatique de textes : Scores cumulatifs}
\framesubtitle{Génération des phrases candidates : Simplification du graphe (Nœuds insignifiants)}

\begin{columns}
	\begin{column}{0.45\textwidth}
		
		\[noeud\_faible(v_i) = \]
		\[( \sum_{(v_i, v_j) \in E} w_{ij} < \frac{1}{MImpN(v_i)} )\]
		
	\end{column}
	\begin{column}{0.45\textwidth}
	  \hgraphpage{img/sim-vertexCut.pdf}
	\end{column}
\end{columns}
	
\end{frame}

\begin{frame}
\frametitle{Résumé automatique de textes : Scores cumulatifs}
\framesubtitle{Génération des phrases candidates : Simplification du graphe (Arcs faibles)}
	
\begin{columns}
	\begin{column}{0.45\textwidth}
		\[arc\_faible(v_i, v_j) = \]
		\[( w_{ij} < \frac{Threshold}{MImpN(v_i)})\]
	\end{column}
	\begin{column}{0.45\textwidth}
		\hgraphpage{img/sim-edgeCut.pdf}
	\end{column}
\end{columns}
		
\end{frame}

\subsection{Notation des phrases} 

\begin{frame}
\frametitle{Résumé automatique de textes : Scores cumulatifs}
\framesubtitle{Notation des phrases : Score statistique des phrases}
	
\begin{columns}
	\begin{column}{0.45\textwidth}
		{\color{darkred}\[Score(s_i/ sim) = sim(s_i, C\backslash s_i)\]}
		{\color{darkgreen}\[Score(s_i/ tfisf) = \sqrt{\sum\limits_{w_{ik} \in s_i} (tfisf(w_{ik}))^2}\]}
		{\color{darkblue}\[Score(s_i/ size) = \frac{1}{|s_i|}\]}
		{\color{darkyellow}\[Score(s_i/ pos) = \max (\frac{1}{i}, \frac{1}{|D| - i + 1})\]}
		\[SSF(s_i/ F) = \prod_{f_i \in F} score(s_i/f_i)\]
		
	\end{column}
	\begin{column}{0.45\textwidth}
		\hgraphpage{img/ssf.pdf}
	\end{column}
\end{columns}

\end{frame}

\begin{frame}
\frametitle{Résumé automatique de textes : Scores cumulatifs}
\framesubtitle{Notation des phrases : Score à base de graphe (GC1)}

\begin{block}{Score GC1}
	\[GC1(s_i) = SSF(s_i) + \sum\limits_{(s_i, s_j) \in E} sim(s_i, s_j) * SSF(s_j)\]
\end{block}
\begin{itemize}
	\item Une phrase peut partager ces informations avec ces voisines
	\item Elle ne peut pas partager la même quantité avec tout le monde 
\end{itemize}
%\begin{columns}
%	\begin{column}{0.45\textwidth}
%		
%	\end{column}
%	\begin{column}{0.45\textwidth}
%		
%	\end{column}
%\end{columns}	
	
\end{frame}

\begin{frame}
\frametitle{Résumé automatique de textes : Scores cumulatifs}
\framesubtitle{Notation des phrases : Score à base de graphe (GC2)}

\begin{block}{Score GC2}
	\[GC2(s_i) = SSF(s_i) + \sum\limits_{(s_i, s_j) \in E} SSF(s_j) - \sum\limits_{(s_i, s_j) \notin E} SSF(s_j)\]
\end{block}
\begin{itemize}
	\item Une phrase peut partager ces informations avec ces voisines
	\item Elle sera récompensé par ces voisines
	\item Elle sera pénalisée par les phrases non voisines  
\end{itemize}
%\begin{columns}
%	\begin{column}{0.45\textwidth}
%		
%	\end{column}
%	\begin{column}{0.45\textwidth}
%		
%	\end{column}
%\end{columns}
	
\end{frame}


\begin{frame}
\frametitle{Résumé automatique de textes : Scores cumulatifs}
\framesubtitle{Notation des phrases : Score à base de graphe (GC3)}

\begin{block}{Score GC3}
	\[GC3(s_i) = SSF(s_i) * (1 + |\{(s_i, s_j) \in E\}|)\]
\end{block}
\begin{itemize}
	\item Une phrase avec plusieurs voisines doit avoir plus de score  
\end{itemize}
%\begin{columns}
%	\begin{column}{0.45\textwidth}
%		
%	\end{column}
%	\begin{column}{0.45\textwidth}
%		
%	\end{column}
%\end{columns}	
	
\end{frame}

\begin{frame}
\frametitle{Résumé automatique de textes : Scores cumulatifs}
\framesubtitle{Notation des phrases : Score à base de graphe (GC4)}
	
\begin{block}{Score GC4}
	\[GC4(s_i) = SSF(s_i) * ( a + \sum\limits_{(s_i, s_j) \in E} sim(s_i, s_j)) \text{ où } a \in \{0, 1\}\]
\end{block}
\begin{itemize}
	\item Une phrase avec plusieurs voisines doit avoir plus de score 
	\item Les voisines n'ont pas la même importance vis-à-vis cette phrase
\end{itemize}
%\begin{columns}
%	\begin{column}{0.45\textwidth}
%		
%	\end{column}
%	\begin{column}{0.45\textwidth}
%		
%	\end{column}
%\end{columns}	

\end{frame}

\subsection{Extraction} 

\begin{frame}
\frametitle{Résumé automatique de textes : Scores cumulatifs}
\framesubtitle{Extraction : e0}

\begin{block}{Extraction selon l'odre du score (e0)}
	\[
	suiv_{e0}  =  \arg\max\limits_i ssfgc(s_i) 
	\text{ où } s_i \in C\backslash S
	\]
\end{block}
\begin{itemize}
	\item On respecte l'ordre selon le score SSF-GC
\end{itemize}
%\begin{columns}
%	\begin{column}{0.45\textwidth}
%		
%	\end{column}
%	\begin{column}{0.45\textwidth}
%		
%	\end{column}
%\end{columns}
	
\end{frame}

\begin{frame}
\frametitle{Résumé automatique de textes : Scores cumulatifs}
\framesubtitle{Extraction : e1}

\begin{block}{Extraction sans redondance (e1)}
	\[suiv_{e1}  =  \arg\max\limits_i ssfgc(s_i)\]
	\[\text{ où } s_i \in C\backslash S \text{ et } simil(s_i, dernier_{e1}) < Th\]
\end{block}
\begin{itemize}
	\item On respecte l'ordre selon le score SSF-GC
	\item On n'ajoute pas les phrases redondantes
\end{itemize}
%\begin{columns}
%	\begin{column}{0.45\textwidth}
%		
%	\end{column}
%	\begin{column}{0.45\textwidth}
%		
%	\end{column}
%\end{columns}
	
\end{frame}

\begin{frame}
\frametitle{Résumé automatique de textes : Scores cumulatifs}
\framesubtitle{Extraction : e2}

\begin{block}{Extraction avec graphe (e2)}
	\[suiv_{e2}  =  \arg\max\limits_i ssfgc(s_i)\]
	\[\text{ où } (dernier_{e2}, s_i) \in E\]
\end{block}
\begin{itemize}
	\item On respecte l'ordre selon le score SSF-GC
	\item On respecte la cohérence (relation entre phrases)
\end{itemize}
%\begin{columns}
%	\begin{column}{0.45\textwidth}
%		
%	\end{column}
%	\begin{column}{0.45\textwidth}
%		
%	\end{column}
%\end{columns}	
	
\end{frame}

\begin{frame}
\frametitle{Résumé automatique de textes : Scores cumulatifs}
\framesubtitle{Extraction : e3}

\begin{block}{Extraction avec graphe et similarité max (e3)}
	\[suiv_{e3}  =  \arg\min\limits_i (iord\ ssfgc(s_i) + iord\ simil(dernier_{e3}, i))\]
	\[\text{ où } (dernier_{e3}, s_i) \in E\]
\end{block}
\begin{itemize}
	\item On respecte l'ordre selon le score SSF-GC
	\item On respecte la cohérence (relation entre phrases)
	\item Les phrases consécutives sont les plus similaires
\end{itemize}
%\begin{columns}
%	\begin{column}{0.45\textwidth}
%		
%	\end{column}
%	\begin{column}{0.45\textwidth}
%		
%	\end{column}
%\end{columns}		
	
\end{frame}

\begin{frame}
\frametitle{Résumé automatique de textes : Scores cumulatifs}
\framesubtitle{Extraction : e4}
	
\begin{block}{Extraction avec graphe et non redondance (e4)}
	\[suiv_{e4}  =  \arg\min\limits_i (iord\ ssfgc(s_i) + ord\ simil(dernier{e4}, i))\]
	\[\text{ où } (dernier_{e4}, s_i) \in E\]
\end{block}
\begin{itemize}
	\item On respecte l'ordre selon le score SSF-GC
	\item On respecte la cohérence (relation entre phrases)
	\item On n'ajoute pas les phrases redondantes
\end{itemize}
%\begin{columns}
%	\begin{column}{0.45\textwidth}
%		
%	\end{column}
%	\begin{column}{0.45\textwidth}
%		
%	\end{column}
%\end{columns}	
	
\end{frame}

\begin{frame}
\frametitle{Résumé automatique de textes : Scores cumulatifs}
\framesubtitle{Extraction : e5}

\begin{block}{Extraction avec graphe et voisines (e5)}
	\[suiv_{e5}  =  \arg\max\limits_i \frac{|\{ s_j | (s_i, s_j) \in E \text{ et } s_j \notin S \}|}{iord\ ssfgc(s_i)}\]
	\[\text{ où } (dernier_{e5}, s_i) \in E\]
\end{block}
\begin{itemize}
	\item On respecte l'ordre selon le score SSF-GC
	\item On respecte la cohérence (relation entre phrases)
	\item On sélectionne la phrase qui peut représenter la majorité des phrases qui ne font pas parti du résumé
\end{itemize}
%\begin{columns}
%	\begin{column}{0.45\textwidth}
%		
%	\end{column}
%	\begin{column}{0.45\textwidth}
%		
%	\end{column}
%\end{columns}		
	
\end{frame}


\subsection{Expérimentation}

\begin{frame}
	\frametitle{Résumé automatique de textes : Scores cumulatifs}
	\framesubtitle{Expérimentation}
	
	\begin{itemize}
		\item multiLing'15 MSS train 
		\begin{itemize}
			\item Comparaison entre le baseline et quelques méthodes connues
			\item Choix du seuil de similarité
			\item Capacité du score cumulatif à améliorer les résumés
			\item Choix des méthodes d'extraction pour chaque variante de score
		\end{itemize}
		\item multiLing'15 MSS test 
		\begin{itemize}
			\item Performance des variantes dans le test
			\item Comparaison de GC avec quelques méthodes (anglais)
		\end{itemize}
	\end{itemize}

\end{frame}

\begin{frame}
\frametitle{Résumé automatique de textes : Scores cumulatifs}
\framesubtitle{Expérimentation : Baseline}

\begin{block}{Comparaison entre le baseline et quelques méthodes connues}
	\scriptsize
	\begin{tabular}{llllllll}
		
		\hline \hline
		
		\multirow{2}{*}{Peer} & \multicolumn{3}{c}{ROUGE-1}  &&  \multicolumn{3}{c}{ROUGE-2}  \\
		\cline{2-4} \cline{6-8} \noalign{\smallskip}
		& R & P & F1 && R & P & F1 \\
		
		\hline
		
		TCC (allsummarizer) & \color{red}0.40863 & \color{red}0.42048 & \color{red}0.41416 && \color{red}0.11529 & \color{red}0.11829 & \color{red}0.11667 \\
		LexRank (linanqiu) & 0.38829 & 0.40238 & 0.39493 && 0.09604 & 0.09908 & 0.09748 \\
		TextRank (summa) & 0.38065 & 0.40803 & 0.39352 && 0.09892 & 0.10548 & 0.10201 \\
		LexRank (sumy) & 0.37926 & 0.39425 & 0.38621 && 0.09350 & 0.09712 & 0.09518 \\
		TextRank (sumy) & 0.37688 & 0.40179 & 0.38864 && 0.09852 & 0.10471 & 0.10145 \\
		SumBasic (sumy) & 0.36909 & 0.37383 & 0.37110 && 0.07683 & 0.07798 & 0.07732 \\
		Random (sumy) & 0.35968 & 0.37472 & 0.36676 && 0.07961 & 0.08262 & 0.08102 \\
		Luhn (sumy) & 0.35671 & 0.39231 & 0.37300 && 0.08575 & 0.09404 & 0.08954 \\
		LSA (sumy) & 0.34674 & 0.37187 & 0.35832 && 0.07678 & 0.08220 & 0.07929 \\
		%KLSummarizer & 0.32745 & 0.34508 & 0.33585 && 0.06317 & 0.06666 & 0.06483 \\
		\hline \hline
	\end{tabular} 
\end{block}
	
\end{frame}

\begin{frame}
\frametitle{Résumé automatique de textes : Scores cumulatifs}
\framesubtitle{Expérimentation : Choix du seuil de similarité}
	
\begin{center}
	\vgraphpage{img/gc-th-choix.pdf}
\end{center}
	
\end{frame}

\begin{frame}
\frametitle{Résumé automatique de textes : Scores cumulatifs}
\framesubtitle{Expérimentation : Capacité du score cumulatif à améliorer les résumés}
	
\begin{center}
	\vgraphpage{img/gc-method-choix.pdf}
\end{center}
	
\end{frame}

\begin{frame}
\frametitle{Résumé automatique de textes : Scores cumulatifs}
\framesubtitle{Expérimentation : Choix des méthodes d'extraction pour chaque variante de score}
	
\begin{block}{Nombre des fois chaque méthode d'extraction maximise ROUGE-1 de chaque GC}
	\centering
	\begin{tabular}{lllllll}
		\hline\hline
		& e0 & e1 & e2 & e3 & e4 & e5 \\
		\hline
		SSF-GC1 & 10 & 4 & 1 & 0 & \textbf{21} & 2 \\
		SSF-GC2 & \textbf{17} & 6 & 2 & 9 & 3 & 1 \\
		SSF-GC3 & 1 & \textbf{15} & 8 & 5 & 4 & 5 \\
		SSF-GC4 & 5 & \textbf{16} & 5 & 4 & 5 & 3 \\
		SSF-GC5 & 4 & \textbf{15} & 3 & 6 & 4 & 6 \\
		\hline\hline
	\end{tabular}  
\end{block}
	
\end{frame}

\begin{frame}
\frametitle{Résumé automatique de textes : Scores cumulatifs}
\framesubtitle{Expérimentation : Performance des variantes dans le test}
	
%\begin{center}
	\hgraphpage{img/gc-test.pdf}
%\end{center}
	
\end{frame}

\begin{frame}
\frametitle{Résumé automatique de textes : Scores cumulatifs}
\framesubtitle{Expérimentation : Comparaison de GC avec quelques méthodes (anglais)}
	
\begin{block}{Comparaison de GC avec quelques méthodes (anglais)}
	\centering\scriptsize
	\begin{tabular}{llllllll}
		\hline\hline
		
		\multirow{2}{*}{Peer} & \multicolumn{3}{c}{ROUGE-1}  &&  \multicolumn{3}{c}{ROUGE-2}  \\
		\cline{2-4} \cline{6-8}\noalign{\smallskip}
		& R & P & F1 && R & P & F1 \\
		
		\hline
		
		LexRank (sumy) & \underline{\color{red}0.41896} & 0.41926 & 0.41879 && 0.11159 & 0.11154 & 0.11148 \\
		SSF-GC1\_mean\_e4 &  0.41635  & \underline{\color{red}0.42816}  & \underline{\color{red}0.42188}  && \underline{\color{red}0.11449}  & \underline{\color{red}0.11764}  & \underline{\color{red}0.11597}\\
		LexRank (linanqiu) & 0.41206 & 0.40646 & 0.40871 && 0.10213 & 0.10087 & 0.10137 \\
		TCC (allsummarizer) & 0.40619 & 0.40866 & 0.40709 && 0.10053 & 0.10115 & 0.10077 \\
		TextRank (summa) & 0.40482 & 0.41731 & 0.41016 && 0.11235 & 0.11515 & 0.11353 \\
		SumBasic (sumy) & 0.39784 & 0.39315 & 0.39518 && 0.09049 & 0.08945 & 0.08989 \\
		TextRank (sumy) & 0.39379 & 0.40310 & 0.39806 && 0.09909 & 0.10087 & 0.09990 \\
		Luhn (sumy) & 0.39014 & 0.40682 & 0.39800 && 0.09478 & 0.09889 & 0.09671 \\
		Random (sumy) & 0.38510 & 0.38874 & 0.38669 && 0.08435 & 0.08514 & 0.08469 \\
		LSA (sumy) & 0.36866 & 0.37727 & 0.37270 && 0.07217 & 0.07393 & 0.07300 \\
		%KLSummarizer & 0.35240 & 0.35701 & 0.35436 && 0.07420 & 0.07524 & 0.07464 \\
		\hline\hline
	\end{tabular} 
\end{block}
	
\end{frame}

% ================================================
%
% ================================================
\section{Test avec l'apprentissage automatique} 

\begin{frame}
	\frametitle{Résumé automatique de textes}
	\framesubtitle{Test avec l'apprentissage automatique}
	
	\begin{center}
		\vgraphpage{img/ml2es-archi.pdf}
	\end{center}
	
\end{frame}

\subsection{Transformation des caractéristiques}

\begin{frame}
\frametitle{Résumé automatique de textes : Test avec ML}
\framesubtitle{Transformation des caractéristiques}

\begin{itemize}
	\item \optword{vecteur vers vecteur}
	\begin{itemize}
		\item coupure des valeurs basses
		\item normalisation min-max 
		{\color{red}
		$\mathbf{minmax(\overrightarrow{v}) = \frac{\overrightarrow{v} - min}{max - min}}$
		}
	\end{itemize}

	\item \optword{vecteur vers scalaire}
	\begin{itemize}
		\item normalisation de sommes phrases/document 
		{\color{red}
		$\mathbf{S2Dsum = \frac{||\overrightarrow{D}|| * sum(\overrightarrow{s})}{||\overrightarrow{s}|| * sum(\overrightarrow{D})}}$
		}
		\item intervalle phrase/document 
		{\color{red}
		$\mathbf{S2Dmxmn = \frac{max(\overrightarrow{s}) - min(\overrightarrow{s})}{max(\overrightarrow{D}) * min(\overrightarrow{D})}}$
		}
		\item normalisation moyenne/max 
		{\color{red}
		$\mathbf{meanMax(\overrightarrow{v}) = \frac{mean(\overrightarrow{v})}{max(\overrightarrow{v})}}$
		}
	\end{itemize}

	\item \optword{scalaire vers scalaire}
	\begin{itemize}
		\item proportion directe 
		{\color{red}
		$\mathbf{DP(x, x_{max}) = (x_{max} - x)/x_{max}}$
		}
		\item proportion inverse {\color{red}$\mathbf{IP(x) = 1/x}$}
		\item séquence géométrique {\color{red}$\mathbf{GS(x) = (0.5)^{x - 1}}$}
		\item normalisation min-max {\color{red}$\mathbf{minmax(x) = \frac{x - min}{max - min}}$}
		\item coupure min et coupure max
	\end{itemize}
\end{itemize}
	
\end{frame}

\subsection{Variations de la méthode}

\begin{frame}
\frametitle{Résumé automatique de textes : Test avec ML}
\framesubtitle{Variations de la méthode : Basique}

\vspace{-6pt}
\begin{center}
	\begin{tabular}{cc}
	\includegraphics[height=.23\textwidth]{img/ml2sa-basic-a.pdf}
	&
	\includegraphics[height=.23\textwidth]{img/ml2sa-basic-b.pdf}
	\\
	Détection Langue & TF ou Similarité\\
	\includegraphics[height=.23\textwidth]{img/ml2sa-basic-c.pdf}
	&
	\includegraphics[height=.23\textwidth]{img/ml2sa-basic-d.pdf}
	\\
	Position & Taille \\
\end{tabular}
\end{center}
	
\end{frame}

\begin{frame}
\frametitle{Résumé automatique de textes : Test avec ML}
\framesubtitle{Variations de la méthode : Filtrage et normalisation}

\begin{itemize}
	\item \optword{Avec filtrage}
	\begin{itemize}
		\item utilisation d'une coupure des valeurs inférieure de la moyenne 
		\item vecteurs concernés : \textcolor{red}{doc\_tf\_seq}, \textcolor{red}{doc\_sim\_seq}, \textcolor{red}{doc\_size\_seq}, \textcolor{red}{sent\_tf\_seq} et \textcolor{red}{sent\_sim\_seq}
	\end{itemize}

	\item \optword{Avec normalisation et filtrage}
	\begin{itemize}
		\item application d'une normalisation min-max ensuite un filtrage sur les vecteurs : \textcolor{red}{doc\_tf\_seq}, \textcolor{red}{doc\_size\_seq} et \textcolor{red}{sent\_tf\_seq}. 
		\item application de filtrage sur : \textcolor{red}{doc\_sim\_seq} et \textcolor{red}{sent\_sim\_seq}. 
		\item normalisation du scalaire \textcolor{red}{sent\_size} où \textit{max} (\textit{min}) est la taille maximale (minimale) des phrases du document.
	\end{itemize}
\end{itemize}
	
\end{frame}

\begin{frame}
\frametitle{Résumé automatique de textes : Test avec ML}
\framesubtitle{Variations de la méthode : Avec caractéristiques scalaires}
	
\hgraphpage{img/ml2es-pure.pdf}

\end{frame}

\subsection{Expérimentation}

\begin{frame}
	\frametitle{Résumé automatique de textes : Test avec ML}
	\framesubtitle{Expérimentation}
	
	\begin{itemize}
		\item Préparer un nouveau corpus prétraité à partir de MultiLing'15 MSS
		\item Comparer les différentes variantes de ML2ExtraSum et TCC
		\item Utiliser les scores intermédiaires pour extraire les phrases
		\item Tester la détection des langues dans chaque variante
	\end{itemize}
	
\end{frame}

\begin{frame}
\frametitle{Résumé automatique de textes : Test avec ML}
\framesubtitle{Expérimentation : Comparaison entre les variantes de ML2ExtraSum et TCC}

\hgraphpage{img/ml2es-test-m.pdf}
	
\end{frame}

\begin{frame}
\frametitle{Résumé automatique de textes : Test avec ML}
\framesubtitle{Expérimentation : Scores intermédiaires}
	
\hgraphpage{img/ml2es-test-s.pdf}
	
\end{frame}


\begin{frame}
\frametitle{Résumé automatique de textes : Test avec ML}
\framesubtitle{Expérimentation : Détection de langues basée sur le RAT (basic)}

\hgraphpage{img/ml2es-lang-basic.png}

\hgraphpage{img/ml2es-lang-legend.png}	

	
\end{frame}

\begin{frame}
\frametitle{Résumé automatique de textes : Test avec ML}
\framesubtitle{Expérimentation : Détection de langues basée sur le RAT (filter)}
	
\hgraphpage{img/ml2es-lang-filter.png}
	
\hgraphpage{img/ml2es-lang-legend.png}
	
\end{frame}

\begin{frame}
\frametitle{Résumé automatique de textes : Test avec ML}
\framesubtitle{Expérimentation : Détection de langues basée sur le RAT (norm)}
	
\hgraphpage{img/ml2es-lang-norm.png}
	
\hgraphpage{img/ml2es-lang-legend.png}
	
\end{frame}

\begin{frame}
\frametitle{Résumé automatique de textes : Test avec ML}
\framesubtitle{Expérimentation : Détection de langues basée sur le RAT (pure)}
	
\hgraphpage{img/ml2es-lang-pure.png}
	
\hgraphpage{img/ml2es-lang-legend.png}
	
\end{frame}


\appendix

\section*{Conclusion}

\begin{frame}
\frametitle{Résumé automatique de textes : Conclusion}
\framesubtitle{Discussion}
	
\begin{itemize}
	\item Les graphes peuvent...
	\begin{itemize}
		\item être utilisés pour filtrer les phrases non importantes
		\item améliorer le score de pertinence
		\item améliorer l'extraction : quantité d'information, cohérence et non redondance
	\end{itemize}
	\item Les caractéristiques de surface peuvent...
	\begin{itemize}
		\item donner de bonnes résultats dans le contexte multilingue
		\item être utilisées pour détecter la langue d'un document
	\end{itemize}
	\item Enfin ...
	\begin{itemize}
		\item la thèse est tellement longue, on sera motivé par le résumé automatique en la lisant
	\end{itemize}
\end{itemize}
	
\end{frame}

\begin{frame}
	\frametitle{Résumé automatique de textes : Conclusion}
	\framesubtitle{Limites}
	
	\begin{itemize}
		\item Considérer les scores des caractéristiques statistiques comme probabilités peut ne pas être c'est la meilleure méthodes pour les combiner
		\item Les caractéristiques de surface sont limitées lorsqu'on veut représenter le sens
		\item L'utilisation de ces caractéristiques avec l'apprentissage automatique peut ne pas améliorer la tâche
		\item La cohérence n'a pas été testée
	\end{itemize}
	
\end{frame}

\begin{frame}
\frametitle{Résumé automatique de textes : Conclusion}
\framesubtitle{Perspectives}
	
\begin{itemize}
	\item En restant dans le contexte multilingue :
	\begin{itemize}
		\item Représenter le sens des phrases (représentation universelle)
		\item Exploiter l'apprentissage automatique 
		\item Utiliser d'autres caractéristiques plus profondes
	\end{itemize}
	\item  Peut-on générer des résumés abstractifs avec moins de ressources linguistiques ?
\end{itemize}
	
\end{frame}

\begin{frame}
\frametitle{Résumé automatique de textes : Conclusion}
\framesubtitle{Contribution : Publications}

\begin{itemize}
	\item \optword{Graph-based Cumulative Score Using statistical features for multilingual automatic text summarization}
	\begin{itemize}
		\item \keyword{Journal :} International Journal of Data Mining, Modelling and Management
		\item \keyword{Date de publication :} pas encore (acceptée)
	\end{itemize}
	\item \optword{AllSummarizer system at MultiLing 2015 : Multilingual single and multi-document summarization}
	\begin{itemize}
		\item \keyword{Conférence :} Special Interest Group on Discourse and Dialogue (SIGDIAL), Prague, Répulique tchèque
		\item \keyword{Date de publication :} 4 septembre 2015
	\end{itemize}
	\item \optword{Exploring Graph Bushy Paths to Improve Statistical Multilingual Automatic Text Summarization}
	\begin{itemize}
		\item \keyword{Conférence :} Computational Intelligence and Its Applications (CIIA), Oran, Algérie
		\item \keyword{Date de publication :} 12 Avril 2018
	\end{itemize}
\end{itemize}
	
\end{frame}

\begin{frame}
\frametitle{Résumé automatique de textes : Conclusion}
\framesubtitle{Contribution : Articles non publiés}
	
\begin{itemize}
	\item \optword{Sentence Object Notation: Multilingual sentence notation based on Wordnet.}
	\begin{itemize}
		\item \keyword{Arxiv :} CoRR abs/1801.00984
		\item \keyword{Date de publication :} 2018
	\end{itemize}
	\item \optword{Automatic text summarization: What has been done and what has to be done.}
	\begin{itemize}
		\item \keyword{Arxiv :} CoRR abs/1904.00688
		\item \keyword{Date de publication :} 2019
	\end{itemize}
\end{itemize}
	
\end{frame}

\begin{frame}
\frametitle{Résumé automatique de textes : Conclusion}
\framesubtitle{Contribution : Ressources}
	

\begin{itemize}
	\item \optword{AllSummarizer} : outil de résumé automatique
	\begin{itemize}
		\item Langage de programmtion : Java
		\item \url{https://github.com/kariminf/allsummarizer}
		\item Contient les méthodes TCC et SSF-GC
	\end{itemize}
	\item \optword{LangPi} : outil de prétraitement
	\begin{itemize}
		\item Langage de programmtion : Java
		\item \url{https://github.com/kariminf/langpi}
		\item Normalisation, segmentation, radicalisation et évaluation ROUGE.
	\end{itemize}
	\item \optword{ML2ExtraSum} : outil de résumé automatique
	\begin{itemize}
		\item Langage de programmtion : Python
		\item \url{https://github.com/kariminf/ML2ExtraSum/}
		\item Résumé extractif en utilisant l'apprentissage automatique + MultiLing’15 MSS corpus
	\end{itemize}
		
\end{itemize}
	
\end{frame}

\begin{frame}[plain]

\begin{center}
	\Huge\bfseries Merci pour votre attention 
\end{center}

\begin{flushright}
	\parbox{.5\textwidth}{%
		Si on a utilisé le résumé automatique, on aura fini la présentation dans 5 minutes
	}%
\end{flushright}

	
\end{frame}

\section*{Bibliographie}
\begin{multicols*}{2}[\frametitle{\insertsection} \usebeamertemplate*{frametitle}]%\usebeamertemplate*{frametitle}\frametitle{Références}
\tiny
\bibliography{cite}
\end{multicols*}

\end{document}

